\section{Introduction}
\label{sec:introduction}

Financial markets are complex adaptive systems. Their observed price dynamics are generated by heterogeneous investors, institutional constraints, liquidity cycles, information flows and feedback mechanisms that operate at several time scales. Even when the empirical object is reduced to a panel of daily returns, the market is not a collection of independent univariate processes. It contains collective motion, sectoral organization, idiosyncratic shocks, crisis amplification and persistent volatility. A central task of empirical econophysics is to identify which part of this structure is robust enough to be interpreted and which part is dominated by noise.

Dependence is particularly important because financial risk is inherently cross-sectional. The volatility of one asset matters, but portfolio risk, contagion, diversification and systemic vulnerability depend on how assets move together. During calm periods, correlations may appear moderate and dispersed. During crises, common movement often increases, and the effective dimensionality of the market can collapse. A method that only studies isolated return series misses this collective component. Conversely, a method that only studies the full sample correlation matrix may confuse persistent structure with sampling noise.

Emerging equity markets make this problem more interesting. They combine global risk channels with local institutional and macroeconomic features. Brazil is a useful case because B3 includes globally exposed commodity companies, large financial institutions, regulated utilities, consumer cyclicals, industrial firms, telecommunications firms and real-estate companies. These assets share exposure to domestic monetary policy, exchange-rate dynamics and country risk, but they also respond to distinct sectoral fundamentals. The resulting dependency matrix is therefore expected to contain a dominant market mode, sectoral modes and a large residual component. Despite these attractive properties, the Brazilian equity market remains underexplored in the complex-systems literature. While studies have addressed Brazilian interbank networks \cite{tabak_2010_topology}, global financial system linkages \cite{silva_2016_structure} and agrarian cross-correlations \cite{stosic_2015_correlations}, a comprehensive econophysics pipeline combining RMT filtering, PMFG construction and volatility forecasting has not been systematically applied to a broad B3 equity universe.

The empirical difficulty is that correlation matrices are noisy. When the number of assets is not negligible relative to the number of observations, a nontrivial part of the eigenspectrum can be compatible with random sampling variation. Random Matrix Theory (RMT) provides a benchmark for this problem by comparing empirical eigenvalues with the spectrum expected from a random correlation matrix \cite{marcenko_1967_distribution,laloux_1999_noise,plerou_2002_random}. Eigenvalues outside the random band are not automatically economic facts, but they identify modes that deserve interpretation. This makes RMT a useful intermediate layer between raw correlations and financial network construction.

Network filtering provides a complementary representation. The Minimum Spanning Tree (MST) converts correlations into a sparse connected backbone with $N-1$ edges \cite{mantegna_1999_hierarchical}. The Planar Maximally Filtered Graph (PMFG) retains a richer planar structure with $3N-6$ edges and preserves triangles and 4-cliques \cite{tumminello_2005_tool}. These networks allow the researcher to ask whether the strongest dependencies are organized by sectors, whether hubs change after market-mode removal and whether the mesoscopic topology contains interpretable information.

This paper uses Brazilian equities to connect three empirical layers. The first layer documents stylized facts, static correlations, sectoral correlation differences and rolling synchronization. The second layer applies RMT to separate market, group/sector and noise components, then studies clustering and filtered matrices. The third layer converts these matrices into MST and PMFG networks and evaluates whether the resulting topological features add information to realized-volatility forecasting.

The forecasting extension is intentionally conservative. Classical volatility models such as GARCH(1,1) \cite{bollerslev_1986_generalized} and HAR-RV \cite{corsi_2009_simple} are strong benchmarks because volatility is persistent. A network-based forecasting model should not be expected to replace this persistence. The more realistic question is whether topology adds incremental information when combined with classical predictors. This paper therefore treats machine learning as an empirical test of the informational content of RMT-filtered networks, not as the main conceptual contribution.

The contribution of this paper is fourfold:
\begin{enumerate}
  \item It builds a reproducible econophysics pipeline for B3 equities using adjusted daily prices from 1998 to 2025, including return construction, correlation diagnostics, sectoral comparison and rolling synchronization.
  \item It applies RMT to a 58-asset historical core universe and separates the correlation matrix into market, group/sector and noise components.
  \item It compares MST and PMFG networks built from original and RMT-filtered matrices, with emphasis on hub reorganization, clique structure, clustering and subsector aggregation.
  \item It evaluates whether RMT- and PMFG-derived features add incremental predictive information to realized-volatility forecasting relative to GARCH(1,1), HAR-RV and classical volatility predictors, using a nested ablation design across three feature sets.
\end{enumerate}

The article is organized as follows. Section~\ref{sec:related_work} reviews financial dependencies, RMT, filtered financial networks and volatility forecasting. Section~\ref{sec:data} describes the data and asset universe. Section~\ref{sec:methodology} presents the empirical pipeline. Sections~\ref{sec:stylized_correlation}--\ref{sec:forecasting} report the results. Sections~\ref{sec:discussion} and~\ref{sec:limitations} discuss interpretation and limitations, and Section~\ref{sec:conclusion} concludes.
